\vspace*{-3ex}
\section{Related Work}
\label{sec:related}
\vspace*{-1.5ex}

Approaches to hierarchical reinforcement learning include learning low-level policies on a collection of easier pre-training tasks \citep{heess2016modulated,tessler2017dsn,frans2017mlsh,rao2021helms,veeriah2021modac}, discovering latent skills via mutual-information maximization \citep{gregor2016vic,florensa2017snn,hausman2018skills,eysenbach2018diayn,achiam2018valor,merel2018humanoid,laskin2022cic,sharma2019dads,xie2020lsp,hafner2020apd,strouse2021disdain}, or training the low-level as a goal-conditioned policy \citep{andrychowicz2017her,levy2017hac,nachum2018hiro,nachum2018norl,co2018sectar,warde2018discern,nair2018rig,pong2019skewfit,hartikainen2019ddl,gehring2021hsd3,shah2021recon}. 
These approaches are described in more detail in \cref{sec:morerelated}.

Relatively few works have demonstrated successful learning of hierarchical behaviors directly from pixels without domain-specific knowledge, such as global XY positions, manually specified pre-training tasks, or precollected diverse experience datasets.
HSD-3 \citep{gehring2021hsd3} showed transfer benefits for low-dimensional control tasks. 
HAC \citep{levy2017hac} learned interpretable hierarchies but required semantic goal spaces. 
FuN \citep{vezhnevets2017fun} learned a two-level policy where both levels maximize task reward and the lower level is regularized by a goal reward but did not demonstrate clear benefits over an LSTM baseline \citep{hochreiter1997lstm}. 
We leverage explicit representation learning and temporally-abstract exploration, demonstrate substantial benefits over flat policies on sparse reward tasks, and underline the generality of our method by showing successful learning without giving task reward to the worker across many domains.
